% Generated draft: requires mathematical and source review.
\documentclass[11pt]{article}
\usepackage[margin=25mm]{geometry}
\usepackage{fontspec}
\setmainfont{Latin Modern Roman}
\usepackage{amsmath,amssymb,mathtools,unicode-math}
\setmathfont{Latin Modern Math}
\usepackage{enumitem,needspace,xurl,hyperref}
\setlength{\parindent}{0pt}
\setlength{\parskip}{5pt}
\setlength{\emergencystretch}{3em}
\newcommand{\N}{\mathbb{N}}
\newcommand{\Z}{\mathbb{Z}}
\newcommand{\Q}{\mathbb{Q}}
\newcommand{\R}{\mathbb{R}}
\newcommand{\C}{\mathbb{C}}
\newcommand{\F}{\mathbb{F}}
\newcommand{\A}{\mathbb{A}}
\newcommand{\PP}{\mathbb{P}}
\newcommand{\E}{\mathbb{E}}
\newcommand{\eps}{\varepsilon}
\newcommand{\dd}{\,\mathrm{d}}
\begin{document}
\section*{UnsolvedMath ID: 10400154}
\textbf{Source problem code:} catalogue entry
\subsection{Definitions and mathematical statement}
Write
\[
 V_k(\Sigma)=H^0\!\left(\mathcal M_\Sigma,L^{\otimes k}\right)
\]
for the space in the source. In the holomorphic interpretation, \(H^0\) is the vector space of global holomorphic sections of the line bundle \(L^{\otimes k}\) on \(\mathcal M_\Sigma\). The record does not define the space \(\mathcal M_\Sigma\), the line bundle \(L\), their dependence on \(\Sigma\), or the allowed levels \(k\); this interpretation is conditional on that omitted geometric setup.

\emph{Problem.} Find a geometric construction of a topological quantum field theory with state spaces \(V_k(\Sigma)\). In particular, give a geometric rule assigning to a \(3\)-manifold \(M\) with boundary identified with \(\Sigma\) a vector
\[
 Z_k(M)\in H^0\!\left(\mathcal M_\Sigma,L^{\otimes k}\right),
 \qquad \partial M=\Sigma,
\]
as part of a TQFT assignment compatible with the relevant equivalences, disjoint union, and gluing of bordisms. The problem asks for a geometric construction, not merely for the choice of some vector in each state space. Any additional boundary, orientation, complex-structure, framing, or anomaly data needed to define these spaces and maps must be specified by a proposed construction; they are not fixed by this excerpt.
\subsection{Short English statement}
Construct the TQFT geometrically by assigning a section of \(L^{\otimes k}\) over \(\mathcal M_\Sigma\) to a \(3\)-manifold bounded by \(\Sigma\), in a way compatible with the TQFT gluing rules. The underlying moduli space, line bundle, and level conventions are missing from the supplied excerpt.
\subsection{Sources}
\begin{itemize}
\item Local UnsolvedMath registry, permanent ID 10400154, source code catalogue entry, Problem 8.8. No individual proposer is supplied.
\item Ohtsuki, \emph{Problems on invariants of knots and 3-manifolds} (2002), Problem 8.8, PDF page 124. Supplied URL: \url{https://msp.org/gtm/2002/04/gtm-2002-04-024s.pdf}; recorded access 2026-07-29; extraction pdf-font-structure. The collection identifies 231 numbered bold headings for problems, conjectures, and questions.
\item Supplied status: open and OPEN-TRIAGE. The research summary describes the geometric construction using \(H^0(\mathcal M_\Sigma,L^{\otimes k})\). Current primary-literature status and redistribution terms for the publicly accessible author/publisher copy remain NEEDS\_REVIEW.
\end{itemize}
\end{document}
